\chapter{Conclusion and Outlook}
\label{ch:conclusion}

\section{Synthesis of the research findings}
\label{sec:assessment_observed_findings}

The central research problem of this dissertation was how rapid activated sludge predictions can be used for engineering analysis when prediction accuracy alone does not establish either a usable treatment state or a good operating decision. The work addressed this problem by connecting four capabilities: complete component prediction, physical enforcement, explicit regression, and independent assessment of selected controls. Together, these capabilities form a common component-state framework, but the results show that each still requires its own evidence. A prediction can be physically consistent yet inaccurate, an inspectable regression can have higher error than a more flexible learner, and a faster operating search can select a worse decision when that decision is evaluated with the original mechanistic model.

The adopted component state retained organic substrate, nitrogen species, phosphorus storage, biomass, and solids, from which COD, TN, TP, and TSS were subsequently obtained through the reporting map \citep{Henze2006}. This ordering allowed the same predicted state to support material accounting, treatment reporting, and operating assessment without treating a reported total as if it uniquely described the underlying component composition. The following sections synthesize the findings in relation to the four research questions in Section~\ref{sec:research_questions}. The corresponding numerical evidence is reported in Sections~\ref{sec:projection_results}, \ref{sec:icsor_results}, and \ref{sec:plant_results}.

\subsection{Complete reactor prediction and data availability}

The first research question examined how accurately different learners reproduced the complete reactor response as the amount of training data changed and as inputs moved beyond the sampled domain. The common reactor benchmark showed that the preferred learner depended on the available simulation budget rather than remaining fixed across design sizes. At 500 cases, XGBoost and LightGBM had the lowest observed raw component nRMSE values of 0.435 and 0.439. At 10,000 cases, the multilayer perceptron instead had the lowest raw and projected values of 0.150 and 0.149. Because these results were obtained from a shared component target rather than separate prediction tasks, they directly address the comparative gap identified in the study. They also show why performance at one training size is insufficient to establish a universally preferred learner \citep{VieringLoog2023,Borisov2024}.

The exterior assessment added a second dimension to this conclusion by examining the usable input range. The multilayer perceptron again had the lowest error, but its projected nRMSE increased to 0.251 under mild exterior inputs and 0.521 under severe exterior inputs. Every projected exterior state nevertheless continued to satisfy the physical checks. The deterioration in severe-domain accuracy therefore arose despite, rather than because of a failure in, physical enforcement. For practical use, this means that a prediction at an unfamiliar loading or operating condition still requires an accuracy assessment even when projection returns a physically admissible state. The benchmark therefore supports model selection within declared data and domain conditions; it does not support unrestricted extrapolation or a single learner ranking that applies everywhere.

\subsection{Physical enforcement and its consequences}

The second research question examined what projection changes when it is applied to the same raw prediction. Every raw assessment state failed the specified mass conservation tolerance, although some predictors returned only non-negative concentrations. This distinction showed that positivity alone was not sufficient to establish physical admissibility. Joint enforcement returned states that passed both requirements in all 766,350 assessed projection instances. These instances included repeated assessment across nested design sizes and therefore did not represent that many distinct simulated reactor states. The maximum invariant residuals were $2.06\times10^{-13}$ in the learning domain and $1.72\times10^{-13}$ in the exterior assessment. The results demonstrate numerical enforcement of the adopted balances and concentration bounds across different learners. They apply the existing projection of \citet{KircherVotsmeier2025} to a common activated sludge prediction task rather than introducing a new projection principle.

Physical compliance, however, did not translate into uniform statistical improvement. At the largest design size, projection improved 188 of 260 predictor--component comparisons, yet aggregate component nRMSE decreased for only five of the thirteen predictors. Across all predictors, TN and TP root errors decreased, whereas COD and TSS root errors increased. Projection therefore closed the physical-compliance gap without eliminating the approximation gap. Its effect also depended on which treatment quantities were of interest: an analysis centered on nutrients could judge the adjustment differently from one centered on organic matter, biomass, or solids. These results support separate reporting of component and composite performance and reinforce the importance of the adjustment metric discussed by \citet{SturmSilva2024}.

Projection displacement and computational cost complete this assessment by showing not only whether a prediction became physically admissible, but also how much it had to change and how much computation was required to return the checked state. Under the stated batch protocol, projection dominated the repeated prediction cost of the fastest raw linear and neural models. A rapid regression evaluation is therefore not equivalent to a rapid complete prediction. The paired assessment provides a basis for selecting a learner and enforcement procedure together while keeping statistical accuracy, physical compliance, adjustment magnitude, and complete prediction cost as distinct criteria \citep{BhosekarIerapetritou2018}.

\subsection{Interpretation of the complete regression response}

The third research question examined how an explicit regression can support interpretable activated sludge prediction while still satisfying mass conservation and non-negativity. ICSOR addressed this problem through named operating, loading, curvature, and interaction terms, followed by learned component coupling and checked projection \citep{Selerio2026}. This staged structure made it possible to trace how an input contribution propagated from the fitted regression to the component state ultimately returned for engineering use. In one fitted model, 44 of 210 unique COD influent-curvature terms exceeded their reporting threshold, while 371 of 380 shared off-diagonal coupling entries exceeded theirs. The fitted representation was therefore selective in influent curvature but dense in coupling. Its interpretive contribution lies in providing an explicit and traceable response structure, not in producing a uniformly sparse model or uniquely identified biological parameters.

The predictive evidence showed where this structure may be practically useful. At 500 cases, ICSOR had the lowest mean held-out pooled composite root error of 10.79, compared with 11.26 for XGBoost and 11.45 for LightGBM. It also had the lowest reported root-error learning-curve area of 6.33 after normalization by the design-size interval. At 10,000 cases, however, the repeated-split errors were 5.94 for ICSOR and 4.62 for the multilayer perceptron, while LightGBM achieved the best mean rank across the broader assessment. These results favor ICSOR when an explicit equation and useful performance under limited simulation data are important, while more flexible learners remain preferable when the principal objective is lower error at the largest data budget. Because accuracy scoring used checked ICSOR states and affine-aligned comparator states, these comparisons do not isolate regression structure from the output-enforcement procedure.

The physical assessment further showed why interpretation must extend beyond the raw regression. Affine projection was sufficient for 436 of the 2000 fixed-partition cases, whereas the remaining 1564 required constrained linear projection. Every final state passed both adopted physical checks. Analytical raw and affine sensitivities, together with the hypothetical active-bound example, showed that a raw regression coefficient does not necessarily represent the final local response returned to the engineer. For example, interpreting an aeration effect requires following the fitted regression through component coupling and then through whichever projection constraints are active at the operating point of interest. ICSOR therefore addresses the interpretation gap through a traceable complete calculation, while questions of coefficient stability and causal process identification remain separate \citep{Selerio2026,Shen2023}.

\subsection{Connected prediction and operating decisions}

The fourth research question extended the problem from an individual reactor to a connected activated sludge plant and asked how joint prediction and projection could support operating decisions. The plant response retained the mixer, five biological reactors, clarifier overflow, clarifier underflow, and stored solids. Joint enforcement connected shared component flows and inventories across these locations rather than checking each unit independently \citep{Alex2008,JeppssonDiehl1996}. The development and holdout designs retained 7386 and 1845 mechanistic states. On the holdout set, projection reduced aggregate nRMSE from 0.04690 to 0.04306, corresponding to an improvement of 8.19\%, and reduced aggregate nMAE by 15.65\%. It also improved 30 of the 32 location--composite nRMSE values, with increases occurring only for COD in the final reactor and overflow. These findings address the network-consistency gap and show that the complete connected procedure produced a broad, although not uniform, statistical benefit.

The operating assessment then revealed an important limitation that average holdout improvement could not resolve. In S4, Extended ICSOR predicted TN and TP concentrations of 2.67 and 2.08 mg L$^{-1}$, whereas independent mechanistic evaluation gave 7.47 and 4.68 mg L$^{-1}$. These discrepancies remained after projection and trust screening. A nutrient requirement lying between the predicted and mechanistically evaluated concentrations would therefore lead to different judgments about whether the selected setting was acceptable. This result directly addresses the decision-quality gap: treatment outcomes must be evaluated at the controls actually selected by the optimization procedure rather than inferred from average holdout performance \citep{Alexandrov1998,Pedrozo2025}.

The common-reference comparison reinforced this conclusion. Smooth NLP produced a lower objective in the nominal case and in all ten influent scenarios. In the nominal case, the objective values were 0.8647 for Extended ICSOR and 0.7862 for Smooth NLP. Extended ICSOR required shorter recorded search times, but those faster searches did not match the objective values of the direct mechanistic route. The answer to the fourth research question is therefore necessarily qualified. Joint prediction and projection can support rapid exploration while maintaining a consistent material account, but they do not remove the need for mechanistic evaluation of the resulting decisions. The comparisons concern available local decisions under the declared treatment and resource priorities; they do not establish a global optimum or universal superiority of either route.

\section{Contributions and gaps addressed}
\label{sec:assessment_conclusions}

The first contribution of the dissertation is a common complete-component reactor benchmark that addresses the lack of directly comparable evidence about learner behavior under changing data availability and input conditions. Rather than assessing only a few reported treatment quantities, the benchmark retains both small component concentrations and the composite quantities derived from them. Its contribution is therefore the defined activated sludge comparison and the evidence it provides about learner behavior, not a claim that one learner is universally preferable.

The second contribution is a paired assessment of a common physical projection across different learners. This assessment makes the distinction between physical enforcement and statistical improvement directly measurable. The projection and stoichiometric foundations themselves come from existing work \citep{SturmWexler2022,KircherVotsmeier2025}; the contribution of this dissertation lies in applying them to complete activated sludge predictions and examining compliance, error redistribution, projection displacement, and computational cost together. In doing so, the study also establishes an important interpretive limit: a physically feasible component state should not be treated as if it were automatically an accurate mechanistic response.

The third contribution is the structured regression and the analysis of its complete returned response. Named terms expose operating and loading relationships, while learned component coupling represents shared dependence among the predicted outputs. Stage-specific physical diagnostics and sensitivities then distinguish what is learned by the regression from what is imposed by projection. This addresses the gap between having visible coefficients and being able to explain the final treatment prediction. The model is therefore inspectable without requiring fitted statistical associations to be interpreted as causal kinetic parameters \citep{Selerio2026}.

The fourth contribution is a connected component-state formulation with joint projection and independent verification of the operating decisions selected from it. The formulation includes shared streams and solids inventories that cannot be represented by a reactor-only output. Independent mechanistic evaluation then connects those predicted states to the treatment and resource consequences of the selected controls. Importantly, the contribution also includes the observed limitations of the surrogate-based decisions. Demonstrating shorter search times alongside higher common-reference objectives provides useful engineering evidence precisely because it avoids assuming that computational speed alone establishes decision superiority.

Table~\ref{tab:integrated_observed_findings} summarizes these gaps, the corresponding findings, and their practical uses. The integration provides a consistent assessment logic across the dissertation, but it does not imply that the numerical scores from the three studies can be combined into a single ranking. Reactor component errors use training-based scales, the structured comparison pools native composite coordinates, and plant holdout errors use location-specific reference ranges. Their validation designs and enforcement procedures also differ. Each set of results must therefore be interpreted within its own assessment rather than treated as directly interchangeable across studies.

\begin{table}[htbp]
\centering\small
\caption{Research gaps, observed answers, and practical uses of the integrated findings.}
\label{tab:integrated_observed_findings}
\begin{tabular}{>{\raggedright\arraybackslash}p{0.24\textwidth}>{\raggedright\arraybackslash}p{0.40\textwidth}>{\raggedright\arraybackslash}p{0.26\textwidth}}
\toprule
Research gap & Observed answer & Practical use\\
\midrule
Uncertain complete-state learner performance & The leading learner changed with data budget, and exterior accuracy deteriorated despite physical compliance & Select for the available data and intended input range\\
Training does not establish physical compliance & All 766,350 projected instances passed both checks, but aggregate error improved for only five of thirteen learners at the largest size & Check balances and bounds separately from treatment accuracy\\
Visible terms do not explain the final response & ICSOR exposed operating and loading terms, while most fixed-partition cases required constrained projection & Interpret the coupled and projected response at the operating point\\
Consistent unit predictions do not establish good plant decisions & Joint prediction improved 30 of 32 local scores, but faster surrogate searches selected higher reference objectives & Verify selected controls with the original plant equations\\
\bottomrule
\end{tabular}
\end{table}

Taken together, these contributions connect process knowledge, statistical approximation, physical enforcement, interpretation, and decision evidence within one assessment framework. Reaction stoichiometry supplies the invariant relations, while unit boundaries provide the transport and inventory relations that connect the treatment system \citep{Henze2006,Alex2008}. Regression approximates the mechanistic response within the sampled conditions, and projection imposes the declared physical requirements on the returned state. Interpretation then follows the complete prediction mapping, while independent mechanistic evaluation establishes the reference treatment outcome at the selected controls. The integrated contribution of the dissertation is this defined connection among the stages and the evidence used to assess each one. It does not claim invention of activated sludge stoichiometry, general projection methods, or surrogate optimization.

\section{Practical implications}
\label{sec:assessment_practical_implications}

\subsection{Treatment and resource priorities}

The operating results illustrate how treatment gains can be considered together with the resources required to obtain them. In S10, the direct decision reduced normalized quality cost from 1.4748 to 0.3699 while increasing the weighted resource contribution from 0.0557 to 0.4191. The total objective consequently decreased from 0.7931 to 0.6040. The selected decision therefore accepted greater capacity and resource contributions in exchange for substantially better weighted treatment. S8 illustrated a different trade-off: a lower resource contribution produced a slightly lower total objective even though all four effluent composites were slightly higher. Presenting the quality and resource contributions separately makes these choices easier to interpret and relate to the treatment priorities represented by the objective \citep{PadronPaez2020,Aboagye2021}.

The selected HRT values provide a corresponding link between operating decisions and capacity planning. Because fresh flow was fixed, varying HRT changed the biological reactor volume. The direct route reached the 36-h upper bound in S5 and S10, whereas the surrogate reached the 6-h lower bound in several cases. Comparing these settings with the resulting treatment performance shows how reactor capacity enters the operating trade-off. Aeration can be interpreted alongside reactor volume because both affect oxygen-transfer capacity, while wasted-solids production can be related to the wasting flow and underflow concentration. At an installed plant, the available reactor volumes and equipment capacities would provide the physical basis for making the same type of comparison among feasible control settings \citep{Henze2008,Miederer2025}.

More generally, separating quality and resource terms provides a practical basis for constructing facility-specific operating objectives. Effluent concentrations and removal values can represent treatment targets, while electricity demand, recycle pumping, sludge handling, and related quantities can represent resource priorities. The assigned weights determine how strongly each quantity influences the preferred operating decision. Presenting the resulting contributions separately allows engineers to explain why one operating condition is preferred and to examine how that preference changes when the treatment or resource priorities change \citep{PadronPaez2020,Miederer2025}.

\subsection{A verified decision workflow}

The results support a workflow in which the surrogate is used primarily for screening and candidate selection rather than as the final authority on the operating decision. The engineer first defines the plant boundary, influent composition, available controls, and treatment targets. The fitted surrogate then generates component responses for comparing alternatives, after which joint projection reconciles the predicted streams and inventories before treatment and resource quantities are evaluated. Physical residuals and projection displacement provide additional diagnostic information that can be used to identify candidates requiring more detailed assessment.

Independent evaluation with the original reaction and settling equations then supplies the final-effluent concentrations, shared solids flows, inventories, and common-reference objective at the selected controls. These quantities place alternative candidates on the same mechanistic basis and therefore provide a consistent foundation for choosing among them. The same evaluations can also identify regions in which additional mechanistic cases would be most useful, allowing subsequent surrogate fitting to concentrate on the operating conditions most relevant to the decisions being made \citep{Alexandrov1998,Hameed2026}.

The component- and location-resolved representation also supports process diagnosis. An increase from fresh-influent to mixer concentrations reveals the material circulated through recycle, while the TSS decrease across the clarifier reflects the contribution of solids separation. Separate ammonium and oxidized-nitrogen concentrations can reveal redistribution among nitrogen forms even when the total nitrogen concentration changes little. Examining these quantities together with returned solids flows and inventories allows engineers to relate aeration, recycle, and wasting decisions to the treatment response rather than relying only on final aggregate indicators \citep{Henze2006,Alex2008,JeppssonDiehl1996}.

\subsection{Computational value}

The measured scenario mean initial-search times were 1.726 s for Extended ICSOR and 5.236 s for Smooth NLP, giving a ratio of approximately 3.03. The shorter surrogate search interval provides computational value when many operating alternatives must be explored, such as across several influent conditions or combinations of aeration, recycle, and wasting settings. In this role, the surrogate reduces the time required for broad screening even though selected candidates still require mechanistic evaluation.

The value of this workflow becomes more relevant when the fitted surrogate is used repeatedly. Model-development effort can then be distributed across many subsequent operating comparisons, while expensive mechanistic calculations are reserved for the candidates most relevant to the decision. The surrogate therefore supports broader exploration, and the mechanistic model supplies detailed verification where it matters. This division of work focuses computational effort on the operating decisions of interest rather than requiring every alternative to be evaluated with the full mechanistic model from the outset \citep{BhosekarIerapetritou2018,Bliek2023}.

\section{Directions for future research}

\subsection{Adaptive sampling around operating decisions}

A natural extension of the present work is to generate additional mechanistic cases in regions where the surrogate most strongly influences the operating decision. An initial surrogate could first screen the declared domain and identify promising control settings. Mechanistic calculations at or near those candidates could then supply additional training states, concentrating new simulations on locations associated with large nutrient discrepancies, strong projection displacements, or treatment requirements that are nearly active. Repeating the fitting and operating-assessment cycle would connect the simulation budget more directly to the decisions the surrogate is intended to support. This direction develops the reference-informed model management of \citet{Alexandrov1998} and the gray-box approximation management discussed by \citet{Hameed2026}.

A useful comparison would hold the mechanistic evaluation budget fixed while contrasting uniform sampling, error-based sampling, and decision-based sampling. The assessment should include effluent errors at the selected controls, the common-reference objective, the number of accepted decisions, and the complete computational time required to deliver them. The nutrient discrepancies observed in S4 provide a concrete case for such a study. Additional mechanistic states around those selected operating conditions could determine whether focused sampling improves TN and TP prediction specifically where those quantities influence the operating choice.

\subsection{Projection aligned with treatment priorities}

The reactor results also motivate further study of the weights used during projection and of the engineering quantities those weights emphasize. A projection that minimizes adjustment in component coordinates can redistribute error differently from one designed around reported nutrients, solids, or other treatment quantities. Future formulations could compare component scaling, composite-sensitive weighting, and weights informed by prediction uncertainty. The underlying mass conservation and non-negativity requirements should remain unchanged so that any differences in treatment accuracy can be attributed to the adjustment metric itself. The species-weighted correction of \citet{SturmSilva2024} provides a precedent for incorporating component magnitude and uncertainty into this choice.

The connected plant provides a broader setting in which to test the same idea. Projection weights could distinguish the importance of internal-state accuracy from that of final-effluent accuracy while still preserving the required unit connections. For example, nutrient-focused weighting could be assessed for its ability to improve predictions at selected controls without introducing large errors in biomass, solids inventory, or other treatment locations. Such a comparison would make the relationship between the physical adjustment and the intended engineering decision more explicit.

\subsection{Uncertainty-aware prediction and operating safeguards}

A further extension is to represent uncertainty alongside the projected point prediction. Repeated fits, model ensembles, or probabilistic component predictions could be used to describe variation in the estimated treatment response, with each sampled state projected onto the same physical set. The resulting distributions could then be evaluated against mechanistic reference concentrations, particularly for TN, TP, and overflow TSS. The conservation-aware probabilistic framework of \citet{Hansen2024} and the differentiable probabilistic projection of \citet{Utkarsh2025} provide starting points for developing such an approach.

These distributions could subsequently be incorporated into operating safeguards. For example, an operating condition might be selected only when an upper predictive bound for TP remains below a specified treatment requirement. Mechanistic and plant-based evaluations could then test how often such safeguards are actually satisfied. The effluent ammonium safeguards considered by \citet{Ren2026} illustrate how treatment requirements can be incorporated into aeration decisions. Separating uncertainty arising from surrogate fitting, influent measurements, and mechanistic parameters would further show which sources of information provide the greatest improvement in decision reliability \citep{Machado2014,Petersen2003}.

\subsection{Stable interpretation of the complete response}

Further work on interpretation can examine whether the fitted coefficient groups remain stable across repeated training samples and different operating domains. The present model already provides natural groups for operating terms, loading terms, curvature, and component coupling. Resampling can reveal which groups remain consistently important and which individual coefficients vary substantially. Comparisons based on the complete predicted response can also identify cases in which different coefficient combinations produce similar treatment behavior. This would directly address the coefficient-stability question identified by \citet{Selerio2026}.

Sensitivity analysis should likewise continue through the final projection rather than stop at the raw regression. Derivatives can be evaluated at representative operating points and checked against direct perturbations of the returned concentrations. Mapping the active constraints would show where a nutrient response changes smoothly and where a projection boundary alters its local slope. Comparing these projected sensitivities with mechanistic sensitivities would provide a stronger connection between model inspection and practical decisions involving aeration, recycle, and wasting \citep{Selerio2026,Shen2023}.

\subsection{Dynamic prediction and control}

The steady-state framework can also be extended to dynamic treatment prediction by replacing fixed component states with trajectories of component concentrations and stored solids. The corresponding governing balances would include the time derivatives of reactor and clarifier inventories \citep{Takacs1991,Alex2008}. Training data could then represent changing influent loads, aeration schedules, recycle settings, and sludge wasting, while physical enforcement would account for the material transported and transformed over each prediction interval.

Such an extension would support predictive-control studies under daily influent variation and other load disturbances. Assessment could include transient nutrient peaks, recovery time, inventory changes, oxygen demand, and the computation required at each control update. The field-tested aeration controller of \citet{Bernardelli2020} and the wastewater control review of \citet{Croll2023} provide relevant examples for developing and evaluating dynamic strategies. The connected component representation developed here provides a useful starting point because it already retains the treatment locations and inventories that influence delayed process responses.

\subsection{Plant validation, topology, and resource models}

Validation with measured plant data is a major next step. Such studies could combine influent fractionation, calibrated mechanistic parameters, concentrations measured at different treatment units, clarifier solids measurements, and recorded operating controls \citep{Petersen2003,Machado2014,Newhart2019}. A prospective assessment could fix the fitting and operating procedures before the final evaluation period is collected \citep{Kaufman2012}. Testing across seasons and loading conditions would then provide evidence about how well the component surrogate transfers from simulated steady-state responses to observed treatment behavior.

The operating objectives can also be adapted to individual facilities. An installed plant has fixed reactor volumes, aeration capacity, recycle capacity, and discharge requirements. Electricity demand, chemical consumption, sludge handling, and operating costs could therefore replace or supplement the engineering proxies used in the present study \citep{Henze2008,Miederer2025}. Multiobjective studies could present the attainable combinations of treatment performance and resource demand before a specific preference is imposed. Extensions to other reactor arrangements, attached-growth systems, and resource-recovery connections could retain the same underlying principle: define the relevant component states and boundary streams before imposing physical enforcement \citep{Henze2008,Aboagye2021,Durkin2024}.

More extensive optimization comparisons could likewise consider several initial control vectors, hybrid surrogate--mechanistic searches, and the fifteen-start basin-sensitivity design described in Chapter~\ref{ch:plant}. Recording the costs of state generation, fitting, prediction, projection, search, and verification would provide a fuller basis for determining when surrogate development becomes worthwhile for repeated decisions \citep{BhosekarIerapetritou2018,Bliek2023}. The central comparison would remain the treatment outcome and common-reference objective obtained for a given computational budget.

\section{Concluding perspective}

This dissertation shows how complete component prediction, physical enforcement, interpretation, and decision verification can be brought together within one activated sludge analysis framework. Its findings answer the research questions without allowing one capability to stand as evidence for another. Learner performance depended on the amount of available data and on the input range being considered. Projection reliably enforced the adopted physical requirements, but it did not improve every prediction error. Structured regression made the treatment response inspectable, but it did not achieve the lowest error at the largest data budget. At the plant scale, connected projection generally improved prediction, while independent mechanistic evaluation showed that the faster surrogate-based searches selected decisions with higher common-reference objectives.

The practical contribution is therefore a defined workflow for rapid exploration with explicit checks on both the predicted state and the operating decision selected from it. A surrogate can be used to screen a broad set of alternatives efficiently, while the original mechanistic model remains the reference for evaluating treatment performance at promising candidate controls. Extending this workflow to actual plant operation will require measured plant data, facility-specific constraints, and locally relevant treatment and resource priorities. Within those limits, computational efficiency has value when it helps deliver an acceptable and verifiable treatment decision, not merely when it reduces search time.